\bookchapter{第四十八章：認識安卓——從作業系統到開發生態}

\chapterintro{恭喜你做出這個決定——從“玩手機的人”變成“做手機應用的人”。在動手之前，咱們先把地圖攤開：安卓從哪裡來？為什麼全世界的開發者都在給它寫應用？一個 App 從你的手指敲鍵盤到使用者手機上，中間都發生了什麼？}

\booksection{一、安卓的身世：從“小機器人”到“全球最大系統”}

安卓（Android）的故事始於 2003 年——一個叫 Andy Rubin 的工程師創辦了 Android 公司，目標很樸素：做一套“更懂相機”的智慧作業系統。後來谷歌在 2005 年收購了它，2007 年聯合幾十家廠商成立開放手機聯盟（OHA），2008 年 9 月，第一款安卓手機 HTC Dream（G1）正式發售。

後面的故事大家都知道了：到 2026 年，安卓是全球市場份額最大的移動作業系統，活躍裝置數以十億計。你的手機、電視盒子、車載大屏、智慧手錶，很可能都在跑安卓。這意味著什麼？意味著你寫的每一行程式碼，都有機會被世界上最多的人用到。

\begin{tcbinfo}{小貼士：安卓的開源基因}
安卓的核心程式碼（AOSP，Android Open Source Project）是開源的——任何廠商都能拿去修改定製。所以小米有 MIUI（現在叫澎湃 OS）、華為有鴻蒙相容層、三星有 One UI。這對開發者是好事：你學的是“安卓開發”，不是“某個品牌的開發”，一套功夫，全網通吃。

\end{tcbinfo}

版本號的進化也很有意思：從 1.5 的紙杯蛋糕（Cupcake）開始，每個大版本都以甜點命名——甜甜圈、閃電泡芙、凍酸奶、薑餅……一直到 2019 年的安卓 10 才改成了純數字。到 2026 年 8 月，最新正式版是 Android 17，而開發時你真正需要關心的“最低支援版本”，主流 App 一般定在 Android 8（API 26）到 Android 10（API 29）之間——支援到幾年前的裝置，你的使用者盤子才夠大。

\begin{center}
\begin{tabularx}{\linewidth}{|X|X|X|}
\hline
\textbf{版本} & \textbf{代號/年份} & \textbf{對開發者重要的事} \tabularnewline
\hline
Android 8 (API 26) & Oreo · 2017 & 通知管道、後臺限制——很多老機型的最低線 \tabularnewline
Android 10 (API 29) & 2019 & 分割區儲存登場，檔案訪問規矩大改 \tabularnewline
Android 12 (API 31) & 2021 & Material You 動態取色，介面觀感大升級 \tabularnewline
Android 13 (API 33) & 2022 & 執行時通知權限，跳出視窗問你"準不準發通知" \tabularnewline
Android 14-15 (API 34-35) & 2023-2024 & 預測性返回、16KB 頁、隱私進一步收緊 \tabularnewline
Android 16-17 (API 36-37) & 2025-2026 & 更快的小版本節奏、AI 能力整合（如 Gemini 擴充套件） \tabularnewline
\hline
\end{tabularx}
\end{center}

\begin{center}\small\color{gray}安卓版本簡史（截至 2026 年 8 月，開發重點關注項）\end{center}

\booksection{二、用什麼語言寫：Java 還是 Kotlin？}

安卓開發歷史上出現過兩種主流語言。Java 是“開國元老”——安卓誕生之初就用它；Kotlin 是 2017 年穀歌欽點的“一等公民”，2019 年起官方明確“Kotlin 優先”（Kotlin-first），到 2026 年，新專案幾乎全部用 Kotlin，舊專案也大多在陸續遷移。

說句實在話：如果你是完全零基礎，直接學 Kotlin，別回頭碰 Java。原因有三：一是 Kotlin 程式碼量比 Java 少 30\%-40\%（同樣的功能，Kotlin 更短更清晰）；二是它天生防“空指標”這種新手噩夢；三是所有現代教學、官方檔案、AI 程式設計工具（還記得咱們的 AI 程式設計書嗎？）都在以 Kotlin 為準。Java 只在維護老專案、讀老程式碼時會遇到。

Java 風格（囉嗦版）

\begin{lstlisting}
public class MainActivity extends Activity \{
@Override
protected void onCreate(Bundle b) \{
super.onCreate(b);
setContentView(R.layout.activity);
\}
\}
\end{lstlisting}

Kotlin 風格（清爽版）

\begin{lstlisting}
class MainActivity : ComponentActivity() \{
override fun onCreate(b: Bundle?) \{
super.onCreate(b)
setContent \{ App() \}
\}
\}
\end{lstlisting}

同樣一件事，Kotlin 少了一堆花括號和樣板。第 50 章咱們會正式學它，現在你只需要記住：安卓開發的“普通話”叫 Kotlin。

\booksection{三、一個 App 是怎麼誕生的？}

很多人以為“開發 App”就是開啟電腦狂寫程式碼。實際上，一個 App 從想法到使用者手機上，大致要經過這樣一條流水線：

想法

設計（介面+功能）

寫程式碼

編譯打包（APK/AAB）

測試（模擬器+真機）

釋出（商店/安裝包）

\begin{center}\small\color{gray}圖 1：一個安卓 App 從想法到上架的全流程\end{center}

流程裡有兩個概念你現在就要認識：

\begin{itemize}
  \item APK：安卓應用的安裝包格式（.apk），連點兩下就能裝到手機上，就像 Windows 的 .exe。平時自己測試、給朋友傳著玩，用它。
  \item AAB：Android App Bundle（.aab），上架 Google Play 用的"進階打包格式"。商店會根據使用者的手機型號動態裁剪，讓每個人下載的包更小。自己裝不了，只能交給商店。
\end{itemize}

另外你還會不斷聽到一個詞：SDK（Software Development Kit，軟體開發工具包）。它就是安卓官方提供的一整套“開發裝備”——包括各版本的 API 庫、除錯工具、模擬器映象等。你寫程式碼時呼叫的 Activity、Button 這些類，都來自 SDK。

\booksection{四、安卓技術棧全景：你的“武器清單”}

正式學習前，先認識一下 2026 年主流安卓專案的標準配置——這張清單你現在看不懂沒關係，全書學完回來看，會豁然開朗：

\begin{center}
\begin{tabularx}{\linewidth}{|X|X|X|}
\hline
\textbf{層次} & \textbf{主流選擇（2026）} & \textbf{幹什麼用的} \tabularnewline
\hline
語言 & Kotlin 2.x & 寫邏輯的語言 \tabularnewline
介面 & Jetpack Compose + Material 3 & 用程式碼畫介面（本書主線） \tabularnewline
架構 & MVVM + ViewModel & 把介面和邏輯分開，程式碼不亂 \tabularnewline
資料儲存 & Room（資料庫）+ DataStore（配置） & 把資料存進手機 \tabularnewline
網路 & Retrofit + OkHttp + Kotlin 協程 & 和伺服器對話 \tabularnewline
構建工具 & Gradle + Kotlin DSL & 把程式碼變成安裝包 \tabularnewline
開發環境 & Android Studio & 谷歌官方的"全家桶"IDE \tabularnewline
\hline
\end{tabularx}
\end{center}

\begin{center}\small\color{gray}2026 年安卓主流技術棧速覽\end{center}

別被這一串英文嚇到。它們就像廚房裡的鍋碗瓢盆——你現在只需要知道“做飯要進廚房”，具體每樣工具怎麼用，後面章節逐個拆給你看。而且好訊息是：Android Studio 建立新專案時會自動幫你把這一整套配置好，你只管炒菜。

\booksection{五、安卓 vs iOS：開發者的第一堂“選邊課”}

你手裡可能正同時有安卓和蘋果裝置。作為開發者，兩邊的體驗差異值得提前瞭解——不是“哪個更好”的站隊，而是“哪邊更適合你起步”的現實帳：

\begin{center}
\begin{tabularx}{\linewidth}{|X|X|X|}
\hline
\textbf{維度} & \textbf{安卓開發} & \textbf{iOS 開發} \tabularnewline
\hline
開發電腦 & Windows / Mac / Linux 都行 & 必須 Mac（Xcode 獨佔） \tabularnewline
開發工具 & Android Studio（免費） & Xcode（免費）+ 上架年費 99 美元 \tabularnewline
上架稽核 & 寬鬆，幾小時\textasciitilde{}幾天 & 嚴格，1-7 天是常態 \tabularnewline
裝機試玩 & APK 隨便發、隨便裝 & 分發受限，測試需簽名/TestFlight \tabularnewline
收入環境 & 廣告/國內管道為主 & App Store 內購付費習慣更好 \tabularnewline
學習成本 & 概念較多（四大元件等） & 生態封閉但體系統一、檔案精良 \tabularnewline
\hline
\end{tabularx}
\end{center}

\begin{center}\small\color{gray}安卓與 iOS 開發路徑對比（2026 年 8 月）\end{center}

說句實在話：如果你的電腦不是 Mac，那這道題根本不用選——安卓是唯一選項。如果你有 Mac，本書第八、九兩篇可以對照著讀，程式設計思路是相通的，學完一邊再學另一邊會快很多。

\begin{tcbpractice}{來動手：看看你手機裡的"安卓"}
開啟手機的 [設定] [關於手機]（或“關於本機”），找到“安卓版本”或“Android 版本”那一欄——連續點它幾下，還能觸發系統的彩蛋小遊戲。記下你的版本號，對照上面的版本表看看它屬於哪一代。你的手機，就是你的第一個“開發測試機”。

\end{tcbpractice}

\begin{tcbsummary}{本章小結}
\begin{itemize}
  \item 安卓 2008 年誕生，2026 年是全球市場份額最大的移動系統，活躍裝置數十億
  \item 最新正式版 Android 17；主流 App 最低支援 Android 8-10，使用者盤子才夠大
  \item 新專案一律用 Kotlin：程式碼少 30\%+、防空指標、官方欽定
  \item APK 自己裝，AAB 上商店；SDK 是官方開發裝備包
  \item 2026 標準技術棧：Kotlin + Compose + MVVM + Room + Retrofit + Gradle + Android Studio
\end{itemize}

\end{tcbsummary}
